\documentclass[aspectratio=169,10pt]{beamer}
\usetheme[numbering=fraction,progressbar=frametitle]{metropolis}
\usepackage{fontspec}
\setsansfont{Latin Modern Sans}
\usepackage{amsmath,amssymb,mathtools,graphicx,booktabs}
\usepackage{polyglossia}
\setdefaultlanguage{english}
\setbeamertemplate{navigation symbols}{}
\setbeamerfont{frametitle}{size=\large}
\setbeamerfont{normal text}{size=\normalsize}
\setlength{\parskip}{.45em}
\hypersetup{pdfauthor={},pdfsubject={Lecture slides},pdfkeywords={},pdftitle={Trigonometric functions}}
\providecommand{\N}{\mathbb{N}}
\providecommand{\Z}{\mathbb{Z}}
\providecommand{\D}{\mathbb{D}}
\providecommand{\Q}{\mathbb{Q}}
\providecommand{\R}{\mathbb{R}}
\providecommand{\C}{\mathbb{C}}
\providecommand{\vect}[1]{\overrightarrow{#1}}
\providecommand{\vi}{\vec{\imath}}
\providecommand{\vj}{\vec{\jmath}}
\providecommand{\vk}{\vec{k}}
\providecommand{\norm}[1]{\left\lVert #1\right\rVert}
\providecommand{\vc}[2]{\begin{pmatrix}#1\\#2\end{pmatrix}}
\providecommand{\deter}[2]{\det\left(#1,#2\right)}
\providecommand{\abs}[1]{\left\lvert #1\right\rvert}
\providecommand{\dd}{\,\mathrm{d}}
\providecommand{\e}{\mathrm{e}}
\providecommand{\Card}{\operatorname{Card}}
\providecommand{\E}{\operatorname{E}}
\providecommand{\V}{\operatorname{V}}
\providecommand{\degre}{^{\circ}}
\providecommand{\un}[1]{\,\mathrm{#1}}
\providecommand{\sci}[2]{#1\times10^{#2}}
\providecommand{\pgcd}{\operatorname{PGCD}}
\providecommand{\ppcm}{\operatorname{PPCM}}
\title{Trigonometric functions}
\subtitle{Lecture and worked examples}
\date{}
\author{}
\begin{document}
\begin{frame}[plain]
\vfill
{\usebeamerfont{title}\inserttitle\par}
\medskip
{\usebeamerfont{subtitle}\insertsubtitle\par}
\medskip
{\usebeamercolor[fg]{progress bar}\rule{\textwidth}{.5pt}}
\vfill
\end{frame}
\begin{frame}{Objectives}
Use parity and periodicity; differentiate sine and cosine; solve equations and inequalities; study a periodic function.
\end{frame}
\begin{frame}{Circle and functions}
On the unit circle, the point located by an angle \(x\) in radians has coordinates \((\cos x,\sin x)\). For every real number \(x\),
\[\cos^2x+\sin^2x=1,\quad -1\le\cos x,\sin x\le1.\]
Cosine is even: \(\cos(-x)=\cos x\). Sine is odd: \(\sin(-x)=-\sin x\). Both functions are \(2\pi\)-periodic.
\[\begin{array}{c|ccccc}x&0&\dfrac{\pi}{6}&\dfrac{\pi}{4}&\dfrac{\pi}{3}&\dfrac{\pi}{2}\\\hline
\cos x&1&\dfrac{\sqrt{3}}{2}&\dfrac{\sqrt{2}}{2}&\dfrac{1}{2}&0\\
\sin x&0&\dfrac{1}{2}&\dfrac{\sqrt{2}}{2}&\dfrac{\sqrt{3}}{2}&1
\end{array}\]
\end{frame}
\begin{frame}{Derivatives and variations}
\[(\sin x)'=\cos x,\qquad(\cos x)'=-\sin x.\]\[
(\sin(ax+b))'=a\cos(ax+b),\quad(\cos(ax+b))'=-a\sin(ax+b).\]
On \([-\pi,\pi]\), cosine increases on \([-\pi,0]\) from \(-1\) to \(1\), then decreases on \([0,\pi]\) from \(1\) to \(-1\). Sine decreases on \([-\pi,-\dfrac{\pi}{2}]\), increases on \([-\dfrac{\pi}{2},\dfrac{\pi}{2}]\), then decreases on \([\dfrac{\pi}{2},\pi]\).
\end{frame}
\begin{frame}{Derivatives and variations}
Read the sign of sine and cosine on the circle, then apply the link between derivative and variations. Periodicity extends the curves; it does not remove the need to study an interval of sufficient length.
\end{frame}
\begin{frame}{Equations and inequalities}
If \(|a|>1\), \(\cos x=a\) and \(\sin x=a\) have no real solution. For a known angle \(\theta\),
\[\cos x=\cos\theta\Longleftrightarrow x=\theta+2k\pi\ \text{or}\ x=-\theta+2k\pi,\]
\[\sin x=\sin\theta\Longleftrightarrow x=\theta+2k\pi\ \text{or}\ x=\pi-\theta+2k\pi,\quad k\in\Z.\]
Then restrict the solutions to the required interval and eliminate duplicates.
\end{frame}
\begin{frame}{Equations and inequalities}
\(\cos x=\dfrac{1}{2}\) gives \(x=\pm\dfrac{\pi}{3}\); \(\sin x=\dfrac{1}{2}\) gives \(x=\dfrac{\pi}{6}\) or \(\dfrac{5\pi}{6}\). \(\cos x\le\dfrac{1}{2}\) gives \([-\pi,-\dfrac{\pi}{3}]\cup[\dfrac{\pi}{3},\pi]\); \(\sin x\ge\dfrac{1}{2}\) gives \([\dfrac{\pi}{6},\dfrac{5\pi}{6}]\).
\end{frame}
\begin{frame}{Equations and inequalities}
To solve \(\cos (2x)=\dfrac{1}{2}\) on \([-\pi,\pi]\), set \(t=2x\in[-2\pi,2\pi]\), solve for \(t\), then divide by \(2\). We obtain \(x\in\{-\dfrac{5\pi}{6},-\dfrac{\pi}{6},\dfrac{\pi}{6},\dfrac{5\pi}{6}\}\).
\end{frame}
\begin{frame}{Equations and inequalities}
For \(f(x)=\sin x+\dfrac12\sin(2x)\) on \([0,\pi]\), \(f'=\cos x+\cos(2x)=2\cos^2x+\cos x-1=(2\cos x-1)(\cos x+1)\). The function increases up to \(\dfrac{\pi}{3}\) then decreases; maximum \(\dfrac{3\sqrt{3}}{4}\).
\end{frame}
\begin{frame}{Equations and inequalities}
For every real number \(x\), \(\sin(2x)=2\sin x\cos x\) and \(\cos(2x)=\cos^2x-\sin^2x=2\cos^2x-1\). These identities are recalled here so that their use in the problems is not left implicit.
\end{frame}
\begin{frame}{Equations and inequalities}
\begin{columns}[T,onlytextwidth]
\begin{column}{.57\textwidth}
The evenness of cosine corresponds to a symmetry about the y-axis; the oddness of sine to a point symmetry about the origin. Periodicity extends each curve by translation of vector \((2\pi,0)\).   The coordinates of the point on the circle are \((\cos x,\sin x)\).
\end{column}
\begin{column}{.4\textwidth}
\centering
\includegraphics[width=\linewidth,height=4.5cm,keepaspectratio]{figures/terminale-14-1-en.pdf}
\end{column}
\end{columns}
\end{frame}
\begin{frame}{Research activity}
\textbf{Height of a point on a wheel - 35 min.}

A wheel turns at constant speed in a vertical plane. Its radius is \(2\) m and its centre is \(3\) m above the ground. A point on the wheel is at its lowest position at time \(t=0\). The wheel makes one revolution in \(8\) s. Its height is modelled by \(h(t)=3-2\cos\left(\dfrac{\pi t}{4}\right)\), where \(t\) is in seconds and \(h(t)\) in metres.
\end{frame}
\begin{frame}{Research activity (2)}
\begin{enumerate}
\setcounter{enumi}{0}\item Check that the model gives the stated initial height, a period of \(8\) s and the expected minimum and maximum heights.

\item Compute \(h' (t)\), then study the variations of \(h\) on \([0;8]\).

\item Determine all the times \(t\in[0;8]\) at which the height is \(4\) m.

\item Determine the total time, during this revolution, for which the point is at least \(4\) m above the ground.

\end{enumerate}
\end{frame}
\end{document}
